\documentclass[Screen16to9,17pt]{foils}
\usepackage{zencurity-slides}

\selectlanguage{english}

\usepackage{tabularx}

% ---------------------------------------------------------------------------
% Academy 2026 - Offensive Security subject
% Day 0 (first half): Introduction to Offensive Security
% Content reused from presentations/pentest/pentest-I-introduction.
% Goals + Plan from the Saint Cyber subject plan. FIXME = genuine gaps.
% ---------------------------------------------------------------------------

\begin{document}

\mytitlepage{Introduction to Offensive Security}
{Saint Cyber -- Cyber Akademiet, 2026}

\slide{Goals for today}
\begin{list2}
\item Explain offensive security as an authorised simulation of adversary tools, tactics and techniques
\item Distinguish vulnerability scanning, vulnerability management, penetration testing and simulated attack exercises
\item Apply the core penetration-testing lifecycle
\item Define scope, objectives, testing conditions and rules of engagement before testing
\item Assess vulnerability severity using CVSS and prioritise remediation
\end{list2}

\slide{Plan for today}
\begin{list2}
\item What offensive security is -- and authorised vs unauthorised
\item The three types: vulnerability management \& scanning, penetration testing, simulated attack exercises
\item CVSS and risk-based prioritisation
\item Planning, scope and rules of engagement
\item Reporting to technical and management audiences
\end{list2}

\slide{Hacker -- cracker}

{\bfseries Short answer -- dont discuss this}

%Det lidt længere svar:\\
Yes, originally there was another meaning to hacker, but the media has perverted it and today, and since early 1990s it has meant breaking into stuff for the public

{\color{red}\hlkbig Today a hacker breaks into systems!}

Reference. Spafford, Cheswick, Garfinkel, Stoll, \ldots
- wrote about this and it was lost

Story is interesting and the old meaning is ALSO used in smaller communities, like hacker spaces full of hackers - doing fun and interesting stuff
\begin{list2}
\item \emph{Cuckoo's Egg: Tracking a Spy Through the Maze of Computer
 Espionage},  Clifford Stoll
\item \emph{Hackers: Heroes of the Computer Revolution},
Steven Levy
\item \emph{Practical Unix and Internet Security},
Simson Garfinkel, Gene Spafford, Alan Schwartz
\end{list2}

\slide{Why even do security testing?}

\begin{list1}
\item Lots of security problems
\item Pentesting may be a requirement from external partners -- example VISA PCI standard
\end{list1}

\begin{list2}
\item Boss asking: should we do a security test?
\item CIO: hmm, okay
\item IT Admins: *sigh* -- I know the security sucks in places!
\item Its not your systems -- dont take the criticism personal, but as an opportunity to get things improved
\vskip 1cm
\item Pentest tools are great resources for doing discovery of assets, evaluating the security of large installations quickly -- in short using pentest tools makes you more efficient!
\end{list2}

\centerline{\Large Many see the benefits after doing a pentest, so try it!}

\slide{The three types of offensive security}
\begin{list2}
\item Vulnerability management \& scanning
\item Penetration testing
\item Simulated attack exercises
\end{list2}
% FIXME Keld: not in pentest-I-introduction -- add a sentence on each
% (objective + limitation), subject plan obj 2
\textbf{FIXME:} one line each on objective and limitation

\slide{Benefits of having a planned security test done}

\begin{quote}
Goal of testing is to reduce risk for the systems and secure the organisation\\ from unexpected loss of data, image and increased costs.
\end{quote}

\begin{list1}
\item Intended audience:
\begin{list2}
\item IT-department and technical personnel
\item Management and board
\item External auditors, government, financial control VISA/PCI, the public
\end{list2}
\item Output from testing:
\begin{list2}
\item Reports with technical content and recommendations
\item Executive summary
\end{list2}
\end{list1}

Goal is not to find a scape goat to blame -- management allocates resources

If security is below in places more resources may be needed.

\slide{Planning a pentest}


\begin{list2}
\item Scope -- what is being tested
\item When is the testing done -- time frame, wall clock time
\item Where are the attacks coming from -- log files will contain the attacks\\
but other attacks from other sources are likely during the attacks, which must be blocked
\item We sometimes go broader than the scope -- perhaps checking the router in front with SNMP or doing a small port 80/tcp scan
\item Agree if Denial of Service is to be tested
\item TL;DR Rules of engagement for the project
\end{list2}

\slide{Testing Agreement and Scope Example}

\begin{list1}
\item Usually the scope could include targets like these:
\begin{list2}
\item 192.168.1.1 -- firewall/router
\item 192.168.1.2 -- mail server
\item 192.168.1.3 -- web server
\item Test to be done from monday 1st until friday 5th
\item Testing done from 192.0.2.0/28
\end{list2}
\item When testing web servers and sites, especially API -- please include hostname, URLs, documentation. If not included some sites and functionality will NOT be tested!
\end{list1}

\slide{Rules of engagement -- rules and ethics for security testing}

\begin{list2}
\item NB: big difference between Denmark and other places!
\item Security consultant must not be the cause of new vulnerabilities due to the testing
\item Security consultant must not install new software on systems without previous agreement
\item Security consultant is not to leave insecure system administrator accounts or settings after testing
\item Security consultant always contact the customer in case of high-risk vulnerabilities
\item It is allowed to peek around in the network -- checking if there might be an insecure development or testing server near by
\item If you meet other security problems outside of the scope we still report them, but perhaps in an appendix
\end{list2}

\centerline{In general be careful of other people networks and systems}

\slide{Selecting systems for testing}

\hlkimage{9cm}{overview-routing-customer-2015.png}

\begin{list2}
\item Routers on the way to critical systems and networks -- especially availability
\item Firewall -- is the environment protected sufficiently, discarding probes
\item Mail servers -- relay testing and also critical data
\item Web servers -- holds data, typically has a lot of functionality
\item Cloud systems, storage systems, anywhere data is saved
\end{list2}

\slide{CVSS -- severity scoring}
% FIXME Keld: pentest-I-introduction has no CVSS slide. Add base/temporal/
% environmental + a worked exercise on a current CVE (subject plan obj 6).
\begin{list2}
\item \textbf{FIXME:} what CVSS is and the metric groups
\item \textbf{FIXME:} exercise -- score a current CVE and discuss likelihood x impact
\end{list2}

\slide{Reporting -- results}

\begin{list1}
\item What is in a pentest report:
\begin{list2}
\item Title, Table of contents, formal report thanks
\item Confidentiality agreement – Write ”Confidential” on each page
\item Executive summary – big companies always want this
\item Information about the scan done, what was it
\item Scope and targets
\item Review of all targets – detailed information and recommendations
\item Conclusion – may be more technical
\item Appendices – various information, Whois info about subnets and prefixes
\end{list2}
\item BTW When delivering a report, it is up to the organisation to decide which recommendations to implement
\end{list1}

Sample report available at: \link{https://github.com/kramse/pentest-report}

\slide{Summary \& references}
\begin{list2}
\item Offensive security is authorised, scoped testing that improves the security posture
\item Agree scope and rules of engagement first; report for technical and management audiences
\item \emph{Web Application Security}, Andrew Hoffman, 2020, ISBN: 9781492053118 (WAS)
\item Sample pentest report: \url{https://github.com/kramse/pentest-report}
\end{list2}

\end{document}
